\section{Operads}
\label{sec:operads}

An \emph{operad} $\OO$ consists of:
\begin{itemize}
\item a set $\OO_n$, whose elements we call \emph{$n$-ary operations},
  for each natural $n \in \naturals$;
\item an identity operation $\id \in \OO_1$; and
\item a \emph{composed} operation $g \compose (f_1, \ldots, f_n) \in
  \OO_{\sum_i n_i}$, for every operation composable pair of operation
  $g \in \OO_n$ and sequence of operations $f_i \in \OO_{n_i}$
\end{itemize}
such that composition is associative and the identity is neutral:
\begin{multline*}
h \compose (g_1(f_{11}, \ldots, f_{1m_1}), \ldots, g_n(f_{n1}, \ldots, f_{nm_n}))
\\\qquad=
(h \compose (g_1, \ldots, g_n))\compose(f_{11}, \ldots, f_{1m_1},
\ldots, f_{n1}, \ldots, f_{nm_n})
\qquad
f \compose (\id, \ldots, \id) = f = \id \compose (f)
\end{multline*}
An operad is then just a $1$-object multi-category $\OO$, by
identifying the $n$-ary operations $\OO_n$ with the $n$-ary
multi-morphisms $\OO(\underbrace{\star, \ldots, \star}_{\text{$n$
    times}}; \star)$.

% \begin{example}\examplelabel{monoid operad}
%   Let $\Monoid$ be the following operad:
%   \begin{itemize}
%   \item Every set of operations $\Monoid_n$ has a unique element
%     $(\bigodot_{i = 1}^n)$.
%   \end{itemize}
%   Since there is only one operation possible at each arity, there is
%   only one choice for identities and compositions, and the operad
%   axioms hold.
% \end{example}
The notion of a model may give more intuition about the definition of
operads.  Let $\OO$ be an operad and $\UU$ a multi-category. A
\emph{model} $\A$ for $\OO$ in $\UU$ consists of:
\begin{itemize}
\item an object $\carrier\A \in \Obj\UU$
\item for each operation $f \in \OO_n$, a multi-morphism
  $\A\sem f : (\overbrace{\carrier\A, \ldots, \carrier \A}^{\text{$n$ times}}) \to \carrier\A$
\end{itemize}
such that composition and identities are preserved:
\[
\A\sem{\id} = \id[\carrier\A] \qquad \A\sem{g \compose (f_1, \ldots,
  f_n)} = \A\sem g \compose (\A\sem{f_1}, \ldots, \A\sem{f_n})
\]
Thus a model of an operad $\OO$ in $\UU$ is just a multi-functor $\A :
\OO \to \UU$.

% \begin{example}\examplelabel{monoid operad models}
%   Recall the multicategory $\pair\Set\bigtimes$ from \exampleref{fin
%     prod multi cat}, i.e., the multicategory induced by choosing the
%   finite cartesian product of sets.  A model for the operad $\Monoid$
%   in $\pair\Set\bigtimes$ is then just a monoid. First, given a monoid
%   $\A = \seq{\carrier \A, (*), e}$, define a model structure over
%   $\carrier\A$ by:
%   \begin{itemize}
%   \item $\A\sem{\bigodot_{i=1}^n}(x_1, \ldots, x_n) := x_1 * \cdots * x_n$
%   \end{itemize}
%   with the usual convention about edge cases, i.e.:
%   \begin{itemize}
%   \item $\A\sem{\bigodot_{i=1}^0}() := e$
%   \item $\A\sem{\bigodot_{i=1}^1}(x) := \A\sem{\id}(x) := x$
%   \item $\A\sem{\bigodot_{i=1}^n}(x_1, \ldots, x_n) := x_1 * \cdots
%     * x_n$ when $n > 1$.
%   \end{itemize}
%   Since $\A$ is associative, we have:
%   \begin{align*}
%   \A&\sem{\bigodot_{i=1}^n\compose(\bigodot_{j=1}^{m_1}, \ldots, \bigodot_{j=1}^{m_n})}(x_{11}, \ldots, x_{1m_1}, \ldots, x_{n1}, \ldots, x_{nm_n})
%   \\&{}=
%   \A\sem{\bigodot_{\ell=1}^{\sum_{i=1}^nm_i}}(x_{11}, \ldots, x_{1m_1}, \ldots, x_{n1}, \ldots, x_{nm_n})
%   \\&{}=
%   x_{11} * \cdots * x_{1m_1}*\cdots * x_{n1} *\cdots* x_{nm_n}
%   \\&{}=
%   (x_{11} * \cdots * x_{1m_1})* \cdots * (x_{n1} *\cdots* x_{nm_n})
%   \\&{}=
%   \A\sem{\bigodot_{i=1}^n}
%   (\A\sem{\bigodot_{j=1}^{m_1}}(x_{11}, \ldots, x_{1m_1})
%   ,\ldots,
%   \A\sem{\bigodot_{j=1}^{m_n}}(x_{n1}, \ldots, x_{nm_n})
%   )
%   \\&{}=
%   \A\sem{\bigodot_{i=1}^n} \compose
%   (\A\sem{\bigodot_{j=1}^{m_1}}
%   ,\ldots,
%   \A\sem{\bigodot_{j=1}^{m_n}}
%   )(x_{11}, \ldots, x_{1m_1}, \ldots, x_{n1}, \ldots, x_{nm_n})
%   \end{align*}
%   Similarly, we have:
%   \[
%   \A\sem{\id \compose \bigodot_{i=1}^n}(x_1, \ldots, x_n) =
%   \A\sem{\bigodot_{i=1}^n}(x_1, \ldots, x_n) =
%   \A\sem{\id} \compose \A\sem{\bigodot_{i=1}^n}(x_1, \ldots, x_n)
%   \]
%   and
%   \begin{align*}
%     \A&\sem{\bigodot_{i=1}^n\compose (\id, \ldots, \id)}(x_1, \ldots, x_n)
%     \\&{}=
%     \A\sem{\bigodot_{i=1}^n}(x_1, \ldots, x_n)
%     \\&{}=
%     x_1 *\cdots *x_n = \id (x_1) *\cdots *(x_n)
%     \\&{}=
%   \A\sem{\bigodot_{i=1}^n}\compose (\A\sem\id, \ldots, \A\sem\id)(x_1, \ldots, x_n)
%   \end{align*}
%   and so $\A$ is a model for $\Monoid$.

%   Conversely, given a model $\A$, define:
%   \begin{itemize}
%   \item $e := \A\sem{\bigodot_{i=1}^0}()$
%   \item $x*y := \A\sem{\bigodot_{i=1}^2}(x,y)$
%   \end{itemize}
%   and observe:
%   \begin{itemize}
%   \item Since $\A\sem{\bigodot_{i=1}^1} = \A\sem{\id} = \id$, we
%     have $\A\sem\id x = x$.
%   \item $e*x = \A\sem{\bigodot_{i=1}^2\compose (\bigodot_{i=1}^0,
%     \id)}(x) = \A\sem{\bigodot_{i=1}^1}(x) = x$
%     and similarly $x*e = x$.
%   \item $(x * y) * z = \A\sem{\bigodot_{i=1}^2(\bigodot_{i=1}^2, \id)}(x, y, z) =
%     \A\sem{\bigodot_{i=1}^2(\id, \bigodot_{i=1}^2)}(x, y, z) = x * (y*z)$.
%   \end{itemize}
%   and so $\A$ is a monoid.

%   Both roundtrips are the identity. When $\A$ is a monoid, then the
%   round-trip monoid is given by:
%   \begin{itemize}
%   \item $e' := \A\sem{\bigodot_{i=1}^0}() = e$;
%   \item $x*'y := \A\sem{\bigodot{i=1}^2}(x,y) = x * y$
%   \end{itemize}
%   Conversely, when $\A$ is a model, then the round-trip model gives:
%   \begin{align*}
%     \A'\sem{\bigodot_{i=1}^n}(x_1, \ldots, x_n)
%     &{}= x_1 * \cdots * x_n
%     \\&{}=
%     \A\sem{\bigodot_{i=1}^2\compose
%       (\id, \bigodot_{i=1}^2\compose
%       (\id, \ldots, \bigodot_{i=1}^2\compose(\id, \id)))}(x_1, \ldots, x_n)
%   \\&{}=
%   \A\sem{\bigodot_{i=1}^n}(x_1, \ldots, x_n)
%   \end{align*}
%   And so the two notions coincide.
% \end{example}
% \begin{example}
%   Given an algebraic theory $\TT$, we define an operad $\OO^{\TT}$:
%   \begin{itemize}
%   \item The $n$-ary operations are the $n$-ary terms $\OO^\TT_n := \TT_n$.
%   \item The identity is given by the projection/variable $\id := x$.
%   \item Composition is given by substitution, taking care of the
%     appropriate weakening:
%     \[
%     x_{11}, \ldots, x_{nm_n} \vdash
%     t\compose (s_1, \ldots, s_n)
%     \definedby
%     t\compose \seq{s_1\compose \seq{x_{11}, \ldots, x_{1m_1}},
%         \ldots, s_n\compose \seq{x_{n1}, \ldots, x_{nm_n}}}
%     \]
%   \end{itemize}
%   The $\Monoid$ operad is quite different from the operad for the
%   algebraic theory $\Monoid$ of monoids. The former has only one
%   operation for each arity, whereas the latter has all possible words
%   over $n$-variables:
%   \[
%   \OO^{\Monoid}_n := \set{x_1, \ldots, x_n \vdash x_{i_1} * \cdots * x_{i_k}
%       \suchthat k \in \naturals, 1 \leq i_1, \ldots, i_k \leq n}
%   \]
%   For example, $\Monoid_1$ only has the identity unary operation, but
%   $\OO^{\Monoid}_1$ also has $e, x^2, \ldots, x^n, \ldots$ as unary
%   operations.
% \end{example}
A \emph{homomorphism} $h : \A[1] \to \A[2]$ between $\OO$-models in
$\UU$ consists of:
\begin{itemize}
\item a morphism $\carrier h : (\carrier{\A[1]}) \to \carrier{\A[2]}$
\end{itemize}
such that:
\begin{itemize}
\item for all $f \in \OO_n$, we have:
  \[
  \carrier h\compose(\A[1]\sem f) = \A[2]\sem f(\carrier h, \ldots, \carrier h)
  \]
\end{itemize}
Identities and composition of homomorphisms are homomorphisms, and so
we have a category $\Mod(\OO, \UU)$ of $\OO$-models in $\UU$.
We also have that a homomorphism of algebras is just a multi-natural
transformation between them as multi-fucntors $\quiver{O-homo as multi-natural}$.

% \begin{example}
%   For every algebraic theory $\TT$ and category $\C$ with chosen
%   products $\prod$, $\Mod(\TT, \C) \isomorphic \Mod(\OO^{\TT},
%   \pair\C\prod)$. Indeed, given a $\TT$-algebra $\A$, define an
%   $\OO^{\TT}$-model on the same carrier by sending each term $t$ to
%   its interpretation $\A\sem{t}$. Conversely, given an
%   $\OO^{\TT}$-model $\A$, we already have interpretations for all the
%   $\TT$-terms, and this interpretation validates the equations. A
%   homomorphism of between $\TT$-algebra is then the same thing as a
%   homomorphism between their corresponding $\OO^{\TT}$-models.
%   In particular the models of $\OO^{\Monoid}$ are the monoids in $\C$:
%   \[
%   \Mod(\OO^{\Monoid}, \pair\C\prod)
%   \isomorphic \Mod(\Monoid, C)
%   \isomorphic \Mon(\C)
%   \]
%   We'll show that we also have an isomorphism $\Mod(\Monoid,
%   \pair\C\prod) \isomorphic \Mon(\C)$. Categorifying the arguments in
%   \exampleref{monoid operad models} shows gives a bijection between
%   the objects of the two categories. The notion of a homomorphisms is
%   maintained across this bijection without changing the underlying
%   morphism, and so we have an isomorphism. Therefore, like with
%   algebraic theories, morita equivalence for operads is a coarser
%   notion than adjoint equivalence (in the $2$-category $\MultiCAT$).
% \end{example}

We'll conclude with a $2$-functor $\Mod(\OO, -) : \MultiCAT \to \CAT$.
\begin{itemize}
\item For every multicategory $\UU$, let $\Mod(\OO, \UU)$ be the
  category of $\OO$-models in $\UU$.
\item For every multi-functor $F : \UU[1] \to \UU[2]$, the
  functor $\Mod(\OO, F) : \Mod(\OO, \UU[1]) \to \Mod(\OO, \UU[2])$
  sends:
  \begin{itemize}
  \item each $\OO$-model $\A$ in $\UU[1]$ to the $\OO$-model:
    \begin{itemize}
    \item $\carrier{F\A} \definedby F\carrier\A$
    \item $(F\A)\sem f : (F\carrier\A, \ldots, F\carrier\A)
      \xto{F(\A\sem f)} F\carrier\A$
    \end{itemize}
  \item each $\OO$-homomorphism $h : \A[1] \to \A[2]$ to the
    $\OO$-homomorphism over $F\carrier h : F\A[1] \to F\A[2]$.
  \end{itemize}
\item For every multi-natural transformation $\quiver{multi-natural
  transformation type}$ we have a natural transformation
  $\quiver{2-cell action of Mod(O,-)}$ given, for every model $\A$, by:
  \[
  \Mod(\OO, \alpha)_{\A} : (F\carrier\A) \xto{\alpha_{\carrier\A}} G\carrier\A
  \]
\end{itemize}
\begin{proposition}
\propositionlabel{Mod(O,-) is a 2-functor}
  This data defines a $2$-functor $\Mod(\OO, -) : \MultiCAT \to \CAT$.
\end{proposition}
\begin{proof}
  To see why the action on $1$-cells is well-defined and functorial:
  \begin{itemize}
    \item For any $\OO$-model $\A$ in $\UU$, view the model as a
      multi-functor $\A : \OO \to \UU$, and note that the model $F\A$
      is given by $F\A : \OO \xto{\A} \UU \xto{F} \UU[2]$. Therefore
      it is a model.
    \item For any $\OO$-model homomorphism $h : \A[1] \to \A[2]$ in
      $\UU$, view it as a multi-natural transformation $\quiver{O-homo
      as multi-natural}$ and note that $\Mod(\OO, F)$ sends it to the
      horizontal compositition $\quiver{O-homo image under 2-functor}$.
      Therefore:
      \begin{itemize}
      \item It is a homomorphism (by type-checking); and
      \item $\Mod(\OO, F)$ is functorial, by the interchange law and
        by preservation of composition:
        \[
        \quiver{O-homo 2-functoriality:composition}\qquad
        \quiver{O-homo 2-functoriality:identity}
        \]
      \end{itemize}
  \end{itemize}

  To see why the action on $2$-cells is well-defined and functorial:
  \begin{itemize}
  \item Consider any multi-natural transormation
    $\quiver{multi-natural transformation type}$. Since the
    multicategory $\OO$ has only one object, $\Mod(\OO,\alpha)_{\A} :
    F \compose \A \star \to G\compose\A \star$ would completely
    determine a multi-natural transformation, and it is the
    multi-natural transformation given by whiskering by $\A$:
    \[
    \quiver{Mod(OO,alpha)_A as a whiskering}
    \]
    It is therefore natural, i.e., each component is an $\OO$-homomorphism
    \[
    \Mod(\OO,\alpha)_{\A} : \Mod(\OO, F)\A = F\A \to G\A = \Mod(\OO, G)\A
    \]
  \item To see that $\Mod(\OO,\alpha)$ is natural, take any
    homomorphism $\quiver{O-homo as multi-natural}$. The naturality
    condition then manifests as the two ways to express the horizontal
    composition at its unique component:
    \[
    \quiver{Mod(OO,alpha) naturality as horizontal composition}
    \]
  \item The preservation of vertical and horizontal composition and
    identities of $2$-cells then follows by taking an $\OO$-model,
    then partitioning these pasting diagrams accordingly:
    \[
    \quiver{Mod(OO,-) preservers vertical composition}
    \qquad
    \quiver{Mod(OO,-) preservers horizontal composition}
    \qquad
    \quiver{Mod(OO,-) preservers identities}
    \]
  \end{itemize}
  We therefore have a $2$-functor.
\end{proof}
